\documentclass[11pt]{letter}

\usepackage[margin=1in]{geometry}
\usepackage[T1]{fontenc}
\usepackage{lmodern}
\usepackage{microtype}
\usepackage[hidelinks]{hyperref}
\hypersetup{
	pdftitle={Cover Letter for Governance Requirements for Coordinating Clinical AI Agents},
	pdfauthor={Alexander J. Ryu}
}

\signature{Alexander J. Ryu\\Corresponding author\\Mayo Clinic\\\href{mailto:Ryu.Alexander@mayo.edu}{Ryu.Alexander@mayo.edu}}
\address{Alexander J. Ryu\\Mayo Clinic\\Rochester, Minnesota, USA}
\date{August 30, 2026}

\begin{document}

\begin{letter}{Editors\\\textit{npj Digital Medicine}}

\opening{Dear Editors,}

Please consider our revised Perspective, ``Governance Requirements for Coordinating Clinical AI
Agents,'' by Bruce Changlong Xu, Ayush Chopra, and Alexander J. Ryu.

Clinical AI systems are moving from isolated recommendations toward actions that intersect across
patients, workflows, institutions, and finite resources. The manuscript argues that safety therefore
cannot be inferred from individual-agent performance alone. It introduces Clinical Infrastructure for
Reliable Coordination (CIRC), a non-exhaustive taxonomy of agent-local, interaction, and institutional
risks linked to independent governance controls and observable evidence.

The revision responds directly to the peer reviews. We replaced the original maturity levels with
independent governance levers; removed the toy simulations, oncology vignette, and population-level
clinical prediction claim; added detailed crosswalks to existing standards and alternative
architectures; specified coordination-layer failures, degraded operation, and operational human
oversight; and added concrete definitions, a structured-message example, and regulatory boundaries.
The manuscript now states explicitly that CIRC is neither a standard nor a validated architecture and
that CliniCARE-Bench illustrates only trace-observable agent-local risk.

We believe the revised manuscript fits the journal's Perspective format because it offers a balanced,
evidence-based, and forward-looking framework intended to stimulate discussion and new experimental
approaches. The manuscript contains a 69-word abstract, 2,961 words of main prose, 29 references, six
tables, two figures, and one illustrative structured-message listing. A detailed point-by-point
response accompanies this submission.

Thank you for considering the revised manuscript.

\closing{Sincerely,}

\end{letter}
\end{document}
